\subsection{Casos de Uso de la Subdefensoría}
\label{subsec:casos_uso_subdefensoria}

La figura \ref{fig:casos_uso_subdefensoria} representa las interacciones de la Subdefensoría, centradas en la supervisión jurídica, validación de resoluciones y la firma electrónica de documentos oficiales.

\vspace{0.5cm}

% ================= CU52 =================
\begingroup
\setlength{\parskip}{0pt}
\begin{longtable}{|p{3.5cm}|p{2cm}|p{2.5cm}|p{7cm}|}
\caption{Especificación de Caso de Uso CU52: Consultar panel de subdefensoría}
\label{tab:cu52} \\
\hline
\textbf{Identificador} & CU52 & \textbf{Nombre} & Consultar panel de subdefensoría \\ \hline
\textbf{Actor Principal} & \multicolumn{3}{p{11.5cm}|}{Subdefensoría} \\ \hline
\textbf{Descripción} & \multicolumn{3}{p{11.5cm}|}{Permite visualizar el listado de expedientes que requieren revisión de jerarquía superior o firma.} \\ \hline
\textbf{Disparador} & \multicolumn{3}{p{11.5cm}|}{El usuario accede al módulo de supervisión.} \\ \hline
\textbf{Precondiciones} & \multicolumn{3}{p{11.5cm}|}{Sesión administrativa iniciada (CU25).} \\ \hline
\textbf{Flujo Principal} & \multicolumn{3}{p{11.5cm}|}{1. El usuario solicita la carga de expedientes.\newline 2. El sistema filtra quejas en estatus de revisión o listas para firma.\newline 3. Se despliega la lista con indicadores de urgencia.} \\ \hline
\textbf{Flujos Alternativos} & \multicolumn{3}{p{11.5cm}|}{Búsqueda por número de folio específico.} \\ \hline
\textbf{Situaciones de error} & \multicolumn{3}{p{11.5cm}|}{Error de respuesta del servidor de datos.} \\ \hline
\textbf{Estado en error} & \multicolumn{3}{p{11.5cm}|}{Mensaje de error y reintento de carga.} \\ \hline
\textbf{Postcondiciones} & \multicolumn{3}{p{11.5cm}|}{Visualización de expedientes pendientes de validar.} \\ \hline
\end{longtable}
\endgroup

\vspace{0.8cm}

% ================= CU53 =================
\begingroup
\setlength{\parskip}{0pt}
\begin{longtable}{|p{3.5cm}|p{2cm}|p{2.5cm}|p{7cm}|}
\caption{Especificación de Caso de Uso CU53: Abrir expediente para revisión}
\label{tab:cu53} \\
\hline
\textbf{Identificador} & CU53 & \textbf{Nombre} & Abrir expediente para revisión \\ \hline
\textbf{Actor Principal} & \multicolumn{3}{p{11.5cm}|}{Subdefensoría} \\ \hline
\textbf{Descripción} & \multicolumn{3}{p{11.5cm}|}{Visualización integral del caso, incluyendo el histórico de análisis del abogado y el analista.} \\ \hline
\textbf{Disparador} & \multicolumn{3}{p{11.5cm}|}{Selección de un folio desde el CU52.} \\ \hline
\textbf{Precondiciones} & \multicolumn{3}{p{11.5cm}|}{Expediente disponible en la bandeja de la Subdefensoría.} \\ \hline
\textbf{Flujo Principal} & \multicolumn{3}{p{11.5cm}|}{1. El sistema recupera el expediente electrónico.\newline 2. Se muestran los proyectos de resolución generados por el área jurídica.} \\ \hline
\textbf{Flujos Alternativos} & \multicolumn{3}{p{11.5cm}|}{N/A.} \\ \hline
\textbf{Situaciones de error} & \multicolumn{3}{p{11.5cm}|}{Incapacidad para cargar documentos PDF adjuntos.} \\ \hline
\textbf{Estado en error} & \multicolumn{3}{p{11.5cm}|}{Aviso de archivo corrupto o no encontrado.} \\ \hline
\textbf{Postcondiciones} & \multicolumn{3}{p{11.5cm}|}{Contenido del expediente expuesto para dictaminación.} \\ \hline
\end{longtable}
\endgroup

\vspace{0.8cm}

% ================= CU54 =================
\begingroup
\setlength{\parskip}{0pt}
\begin{longtable}{|p{3.5cm}|p{2cm}|p{2.5cm}|p{7cm}|}
\caption{Especificación de Caso de Uso CU54: Validar proyecto de resolución}
\label{tab:cu54} \\
\hline
\textbf{Identificador} & CU54 & \textbf{Nombre} & Validar proyecto de resolución \\ \hline
\textbf{Actor Principal} & \multicolumn{3}{p{11.5cm}|}{Subdefensoría} \\ \hline
\textbf{Descripción} & \multicolumn{3}{p{11.5cm}|}{Revisión de la legalidad y congruencia del documento de resolución antes de su emisión final.} \\ \hline
\textbf{Disparador} & \multicolumn{3}{p{11.5cm}|}{El usuario inicia el proceso de validación en el expediente abierto.} \\ \hline
\textbf{Precondiciones} & \multicolumn{3}{p{11.5cm}|}{Proyecto de resolución cargado por un abogado.} \\ \hline
\textbf{Flujo Principal} & \multicolumn{3}{p{11.5cm}|}{1. El Subdefensor revisa el texto del proyecto.\newline 2. El usuario decide entre: Aprobar, Rechazar o Solicitar Corrección.} \\ \hline
\textbf{Flujos Alternativos} & \multicolumn{3}{p{11.5cm}|}{\textbf{1a. Correcciones:} Se devuelve al abogado con comentarios.} \\ \hline
\textbf{Situaciones de error} & \multicolumn{3}{p{11.5cm}|}{Fallo al guardar las observaciones de revisión.} \\ \hline
\textbf{Estado en error} & \multicolumn{3}{p{11.5cm}|}{Se solicita al usuario no cerrar la ventana para evitar pérdida de datos.} \\ \hline
\textbf{Postcondiciones} & \multicolumn{3}{p{11.5cm}|}{Estado de validación registrado.} \\ \hline
\end{longtable}
\endgroup

\vspace{0.8cm}

% ================= CU55 =================
\begingroup
\setlength{\parskip}{0pt}
\begin{longtable}{|p{3.5cm}|p{2cm}|p{2.5cm}|p{7cm}|}
\caption{Especificación de Caso de Uso CU55: Firmar resolución electrónicamente}
\label{tab:cu55} \\
\hline
\textbf{Identificador} & CU55 & \textbf{Nombre} & Firmar resolución electrónicamente \\ \hline
\textbf{Actor Principal} & \multicolumn{3}{p{11.5cm}|}{Subdefensoría} \\ \hline
\textbf{Descripción} & \multicolumn{3}{p{11.5cm}|}{Aplicación de la firma electrónica institucional para dar validez legal al documento.} \\ \hline
\textbf{Disparador} & \multicolumn{3}{p{11.5cm}|}{El usuario selecciona la opción "Firmar Documento".} \\ \hline
\textbf{Precondiciones} & \multicolumn{3}{p{11.5cm}|}{Resolución validada previamente y certificado digital vigente.} \\ \hline
\textbf{Flujo Principal} & \multicolumn{3}{p{11.5cm}|}{1. El sistema solicita la clave privada o token.\newline 2. El sistema incluye CU56 (Generar cadena de firma).\newline 3. Se estampa la firma en el PDF.} \\ \hline
\textbf{Flujos Alternativos} & \multicolumn{3}{p{11.5cm}|}{N/A.} \\ \hline
\textbf{Situaciones de error} & \multicolumn{3}{p{11.5cm}|}{Certificado digital caducado; contraseña incorrecta.} \\ \hline
\textbf{Estado en error} & \multicolumn{3}{p{11.5cm}|}{Operación cancelada y mensaje de error de autenticidad.} \\ \hline
\textbf{Postcondiciones} & \textbf{CU57} (Estatus Firmado) y \textbf{CU58} (Notificar). \\ \hline
\end{longtable}
\endgroup

\vspace{0.8cm}

% ================= CU56 =================
\begingroup
\setlength{\parskip}{0pt}
\begin{longtable}{|p{3.5cm}|p{2cm}|p{2.5cm}|p{7cm}|}
\caption{Especificación de Caso de Uso CU56: Generar cadena original y sello}
\label{tab:cu56} \\
\hline
\textbf{Identificador} & CU56 & \textbf{Nombre} & Generar cadena original y sello \\ \hline
\textbf{Actor Principal} & \multicolumn{3}{p{11.5cm}|}{Sistema} \\ \hline
\textbf{Descripción} & \multicolumn{3}{p{11.5cm}|}{Cifrado de los datos del documento para garantizar su integridad (no alteración).} \\ \hline
\textbf{Disparador} & \multicolumn{3}{p{11.5cm}|}{Ejecución de la firma electrónica (CU55).} \\ \hline
\textbf{Precondiciones} & \multicolumn{3}{p{11.5cm}|}{Documento PDF final generado.} \\ \hline
\textbf{Flujo Principal} & \multicolumn{3}{p{11.5cm}|}{1. El sistema extrae el hash del documento.\newline 2. Genera la cadena original según el estándar institucional.\newline 3. Crea el código QR de validación.} \\ \hline
\textbf{Flujos Alternativos} & \multicolumn{3}{p{11.5cm}|}{N/A.} \\ \hline
\textbf{Situaciones de error} & \multicolumn{3}{p{11.5cm}|}{Error en el algoritmo de cifrado.} \\ \hline
\textbf{Estado en error} & \multicolumn{3}{p{11.5cm}|}{Fallo crítico, el documento no se marca como firmado.} \\ \hline
\textbf{Postcondiciones} & \multicolumn{3}{p{11.5cm}|}{Sello digital generado satisfactoriamente.} \\ \hline
\end{longtable}
\endgroup

\vspace{0.8cm}

% ================= CU57 =================
\begingroup
\setlength{\parskip}{0pt}
\begin{longtable}{|p{3.5cm}|p{2cm}|p{2.5cm}|p{7cm}|}
\caption{Especificación de Caso de Uso CU57: Cambiar estatus a firmado}
\label{tab:cu57} \\
\hline
\textbf{Identificador} & CU57 & \textbf{Nombre} & Cambiar estatus a firmado \\ \hline
\textbf{Actor Principal} & \multicolumn{3}{p{11.5cm}|}{Sistema} \\ \hline
\textbf{Descripción} & \multicolumn{3}{p{11.5cm}|}{Actualización del estado de la queja para permitir su notificación final.} \\ \hline
\textbf{Disparador} & \multicolumn{3}{p{11.5cm}|}{Finalización exitosa del CU55.} \\ \hline
\textbf{Precondiciones} & \multicolumn{3}{p{11.5cm}|}{Documento PDF actualizado con sellos digitales.} \\ \hline
\textbf{Flujo Principal} & \multicolumn{3}{p{11.5cm}|}{1. El sistema actualiza el campo estatus a 'FIRMADO Y NOTIFICABLE'.\newline 2. Registra el evento en la bitácora del folio.} \\ \hline
\textbf{Flujos Alternativos} & \multicolumn{3}{p{11.5cm}|}{N/A.} \\ \hline
\textbf{Situaciones de error} & \multicolumn{3}{p{11.5cm}|}{Base de datos ocupada o sin conexión.} \\ \hline
\textbf{Estado en error} & \multicolumn{3}{p{11.5cm}|}{Reintento automático por parte del microservicio.} \\ \hline
\textbf{Postcondiciones} & \multicolumn{3}{p{11.5cm}|}{El expediente queda listo para envío al quejoso.} \\ \hline
\end{longtable}
\endgroup

\vspace{0.8cm}

% ================= CU58 =================
\begingroup
\setlength{\parskip}{0pt}
\begin{longtable}{|p{3.5cm}|p{2cm}|p{2.5cm}|p{7cm}|}
\caption{Especificación de Caso de Uso CU58: Notificar resolución firmada}
\label{tab:cu58} \\
\hline
\textbf{Identificador} & CU58 & \textbf{Nombre} & Notificar resolución firmada \\ \hline
\textbf{Actor Principal} & \multicolumn{3}{p{11.5cm}|}{Sistema} \\ \hline
\textbf{Descripción} & \multicolumn{3}{p{11.5cm}|}{Envío automático del documento final al correo electrónico del quejoso.} \\ \hline
\textbf{Disparador} & \multicolumn{3}{p{11.5cm}|}{Cambio de estatus en el CU57.} \\ \hline
\textbf{Precondiciones} & \multicolumn{3}{p{11.5cm}|}{Correo del destinatario registrado.} \\ \hline
\textbf{Flujo Principal} & \multicolumn{3}{p{11.5cm}|}{1. El sistema adjunta la resolución firmada.\newline 2. Envía la notificación con las instrucciones de descarga.} \\ \hline
\textbf{Flujos Alternativos} & \multicolumn{3}{p{11.5cm}|}{N/A.} \\ \hline
\textbf{Situaciones de error} & \multicolumn{3}{p{11.5cm}|}{Fallo del servidor de correo.} \\ \hline
\textbf{Estado en error} & \multicolumn{3}{p{11.5cm}|}{El sistema marca la notificación como 'Pendiente de reenvío'.} \\ \hline
\textbf{Postcondiciones} & \multicolumn{3}{p{11.5cm}|}{El usuario recibe el resultado oficial de su queja.} \\ \hline
\end{longtable}
\endgroup

\vspace{0.8cm}

% ================= CU59 =================
\begingroup
\setlength{\parskip}{0pt}
\begin{longtable}{|p{3.5cm}|p{2cm}|p{2.5cm}|p{7cm}|}
\caption{Especificación de Caso de Uso CU59: Generar reportes estadísticos}
\label{tab:cu59} \\
\hline
\textbf{Identificador} & CU59 & \textbf{Nombre} & Generar reportes estadísticos \\ \hline
\textbf{Actor Principal} & \multicolumn{3}{p{11.5cm}|}{Subdefensoría} \\ \hline
\textbf{Descripción} & \multicolumn{3}{p{11.5cm}|}{Obtención de datos agregados sobre el desempeño institucional y tipos de quejas.} \\ \hline
\textbf{Disparador} & \multicolumn{3}{p{11.5cm}|}{El usuario selecciona el menú de 'Reportes'.} \\ \hline
\textbf{Precondiciones} & \multicolumn{3}{p{11.5cm}|}{Permisos de nivel directivo.} \\ \hline
\textbf{Flujo Principal} & \multicolumn{3}{p{11.5cm}|}{1. El sistema ofrece filtros (fechas, áreas, estatus).\newline 2. El sistema incluye CU60 (Procesar datos).\newline 3. Se genera archivo descargable.} \\ \hline
\textbf{Flujos Alternativos} & \multicolumn{3}{p{11.5cm}|}{Visualización de gráficas en pantalla.} \\ \hline
\textbf{Situaciones de error} & \multicolumn{3}{p{11.5cm}|}{Timeout por volumen excesivo de datos.} \\ \hline
\textbf{Estado en error} & \multicolumn{3}{p{11.5cm}|}{Solicitud de reducir el rango de fechas.} \\ \hline
\textbf{Postcondiciones} & \multicolumn{3}{p{11.5cm}|}{Reporte generado en formato PDF/Excel.} \\ \hline
\end{longtable}
\endgroup

\vspace{0.8cm}

% ================= CU61 =================
\begingroup
\setlength{\parskip}{0pt}
\begin{longtable}{|p{3.5cm}|p{2cm}|p{2.5cm}|p{7cm}|}
\caption{Especificación de Caso de Uso CU61: Gestionar usuarios administrativos}
\label{tab:cu61} \\
\hline
\textbf{Identificador} & CU61 & \textbf{Nombre} & Gestionar usuarios administrativos \\ \hline
\textbf{Actor Principal} & \multicolumn{3}{p{11.5cm}|}{Subdefensoría} \\ \hline
\textbf{Descripción} & \multicolumn{3}{p{11.5cm}|}{Administración de las cuentas del personal (recepcionistas, analistas, abogados).} \\ \hline
\textbf{Disparador} & \multicolumn{3}{p{11.5cm}|}{Necesidad de alta, baja o cambio de rol de un empleado.} \\ \hline
\textbf{Precondiciones} & \multicolumn{3}{p{11.5cm}|}{Sesión con privilegios de administrador.} \\ \hline
\textbf{Flujo Principal} & \multicolumn{3}{p{11.5cm}|}{1. El usuario accede al padrón de personal.\newline 2. El sistema permite: Crear (CU62), Editar (CU63) o Eliminar (CU64).} \\ \hline
\textbf{Flujos Alternativos} & \multicolumn{3}{p{11.5cm}|}{Búsqueda de usuario por nombre o clave de empleado.} \\ \hline
\textbf{Situaciones de error} & \multicolumn{3}{p{11.5cm}|}{Intento de eliminar una cuenta con expedientes activos asignados.} \\ \hline
\textbf{Estado en error} & \multicolumn{3}{p{11.5cm}|}{Bloqueo de la acción y solicitud de reasignación de quejas.} \\ \hline
\textbf{Postcondiciones} & \multicolumn{3}{p{11.5cm}|}{Base de datos de usuarios actualizada.} \\ \hline
\end{longtable}
\endgroup

\vspace{0.8cm}

% ================= CU62 =================
\begingroup
\setlength{\parskip}{0pt}
\begin{longtable}{|p{3.5cm}|p{2cm}|p{2.5cm}|p{7cm}|}
\caption{Especificación de Caso de Uso CU62: Crear nuevo usuario}
\label{tab:cu62} \\
\hline
\textbf{Identificador} & CU62 & \textbf{Nombre} & Crear nuevo usuario \\ \hline
\textbf{Actor Principal} & \multicolumn{3}{p{11.5cm}|}{Subdefensoría} \\ \hline
\textbf{Descripción} & \multicolumn{3}{p{11.5cm}|}{Registro de un nuevo integrante del equipo administrativo.} \\ \hline
\textbf{Disparador} & \multicolumn{3}{p{11.5cm}|}{Selección de "Nuevo Usuario" en CU61.} \\ \hline
\textbf{Precondiciones} & \multicolumn{3}{p{11.5cm}|}{Disponer del correo institucional del nuevo usuario.} \\ \hline
\textbf{Flujo Principal} & \multicolumn{3}{p{11.5cm}|}{1. Ingreso de datos personales y asignación de rol.\newline 2. El sistema envía invitación de activación por correo.} \\ \hline
\textbf{Flujos Alternativos} & \multicolumn{3}{p{11.5cm}|}{N/A.} \\ \hline
\textbf{Situaciones de error} & \multicolumn{3}{p{11.5cm}|}{Correo electrónico ya registrado en el sistema.} \\ \hline
\textbf{Estado en error} & \multicolumn{3}{p{11.5cm}|}{Mensaje de "Usuario ya existe".} \\ \hline
\textbf{Postcondiciones} & \multicolumn{3}{p{11.5cm}|}{Nuevo registro creado en estatus 'Pendiente de activar'.} \\ \hline
\end{longtable}
\endgroup

\vspace{0.8cm}

% ================= CU39 (Reutilizado) =================
